\documentclass[12pt]{article}

\usepackage[margin=2.5cm]{geometry}

\usepackage[colorlinks,linkcolor=black,citecolor=blue,urlcolor=black]{hyperref}

\usepackage{amsfonts,hyphenat,booktabs,setspace,amssymb,float}

\usepackage{pdfpages}

\usepackage{amsmath, amssymb, mathrsfs}

\usepackage{graphicx, psfrag, epsf}

\usepackage{enumerate}

\usepackage{natbib,longtable}

\usepackage{url} 

\usepackage{subfigure}



\newenvironment{psmallmatrix}

{\left(\begin{smallmatrix}}

{\end{smallmatrix}\right)}

\newcommand{\blind}{0}



\setlength{\bibsep}{0pt}

\def\spacingset#1{\renewcommand{\baselinestretch}%

{#1}\small\normalsize} \spacingset{1.5}



\newcommand{\btheta}{\mbox{\boldmath $\theta$}}

\newcommand{\bmu}{\mbox{\boldmath $\mu$}}

\newcommand{\balpha}{\mbox{\boldmath $\alpha$}}

\newcommand{\bgamma}{\mbox{\boldmath $\gamma$}}

\newcommand{\bbeta}{\mbox{\boldmath $\beta$}}

\newcommand{\bpsi}{\mbox{\boldmath $\psi$}}

\newcommand{\bphi}{\mbox{\boldmath $\phi$}}

\newcommand{\bxi}{\mbox{\boldmath $\xi$}}

\newcommand{\bnu}{\mbox{\boldmath $\nu$}}

\newcommand{\bsigma}{\mbox{\boldmath $\sigma$}} 



\begin{document}



\begin{center}

 {\Large{\bf Joint Mean-Variance Inference for Nonlinear Mixed-Effects Models with 

   Measurement Errors and Outliers}}



      March, 2022 



      Ye et al. 

\end{center} 

\vspace{1cm}



\section{Introduction}



Nonlinear mixed-effects (NLME) models are widely used in many biological applications, 

such as pharmacokinetic analysis, HIV viral dynamics, and studies of growth and decay. 

Many NLME models are derived based on the underlying data-generation mechanisms, so they 

are also called {\em scientific models} or {\em mechanistic models}. They may be more desirable 

than empirical models such as linear mixed effects (LME) models and nonparametric mixed 

models, since they may provide more reliable predicted values for (unobserved) missing/censored values 

or future values. In many longitudinal studies, the within-individual repeated measurements exhibit 

large variations and these variations may change over time. It is important to understand 

the nature of the within-individual systematic and random variations in order to provide 

good predictions and efficient statistical inference. 





Figure 1 shows the profiles of HIV viral load (HIV-1 RNA copies) longitudinal measurements 

on a sub-sample of individuals from an AIDS study during an 

anti-HIV treatment. We see that viral loads decline

rapidly in the initial period (first few weeks), roughly following an exponential pattern. 

The variations of the within-individual repeated measurements 

appear to {\em change} over time. 

Moreover, there seem some {\em outliers} in the data, i.e., some data do not seem 

to follow the general patterns of the whole data. It is interesting to see if these 

large and time-varying within-individual variations may be partially explained by 

covariates such as CD4 cell counts. 

Based on some biological arguments, as well as a set of differential equations which

describe the viral production and elimination processes, Wu and Ding (1999)

showed that the viral load trajectories during the treatment may well be modelled by

an NLME model.





\begin{figure}[h]

    \centering

    \includegraphics[scale=0.6]{/home/lang/Figures/ch9outlier.pdf}

    \caption{\it\small Viral loads trajectories (in log$_{10}$ scale) during 

an anti-HIV treatment.} 

    \label{fig1}

\end{figure}



   To see if the large and time-varying within-individual variations in longitudinal 

viral load data may be partially explained by time-varying 

covariates such as CD4 cell counts, we may model the within-individual variations with 

time-varying covariates and incorporate the large within-individual variations through 

random effects in the variance model. In other words, we may jointly model the mean and variance  

of the longitudinal viral load data: we model the mean based on an NLME model and model the variance 

using another mixed effects model. The advantages of modelling the 

within-individual variations, compared to NLME models without modelling the variance,

 are as follows: (a) we may gain good understanding of 

the nature of the within-individual variation  and 

if this variation may be partially explained by covariates; (b) we may obtain 

more efficient estimates of the main parameters of interest; (c) as will be described 

in Section 2, modeling variance also allows us to conduct robust inference to 

address outliers in the data; and (d) we may provide 

more accurate predictions for the missing/censored response values or future values. 



   For example, by modelling the variance, we may obtain an more efficient 

estimate of initial viral decay rates, which allows us to better evaluate and compare 

different treatments since the initial viral decay rates may reflect the 

efficacies of the treatments (Wu and Ding, 1999). Understanding the nature of the 

within-individual variations in viral load also allows us to choose good timing for 

treatment interruption. 



    There is some literature on joint modelling mean and variance for longitudinal 

data (e.g., Lin et al. 1997; Pourahmadi  1999; Hedeker et al. 2008; German et al. 2022). 

To our knowledge, all previous work focused on {\em linear} models, either with random effects 

or without random effects. In this article, we extend previous work to NLME models, and 

we also address covariate measurement errors and outlying responses. NLME models are 

{\em nonlinear} in the random effects and the dimension of the 

random effects may be high, so computation for the joint model can be 

challenging. In this article, we propose a computationally efficient approximate 

method for approximate likelihood 

inference. 



    The article is organized as follows ...







\section{The Models}



\subsection{The mean and variance models}



      Let $y_{ij}$ be the response value for individual $i$ at time $t_{ij}$, and let 

 ${\bf y}_i = (y_{i1}, y_{i2}, \cdots, y_{in_i})^T$,  $i=1,2,\cdots,n; \; j=1,2,\cdots,n_i$.  

We consider the following 

nonlinear mixed effects (NLME) model

\begin{eqnarray}

     {\bf y}_i = g({\bf x}_i, \bbeta, {\bf u}_i) + {\bf e}_i, \qquad\qquad i=1,2,\cdots,n; \;

j=1,2,\cdots,n_i,   \label{nlme}

\end{eqnarray}

where $g(\cdot)$ is a known nonlinear function for the mean $E({\bf y}_i)$, ${\bf x}_i$ contains covariates including 

time, vector $\bbeta = (\beta_1,\beta_2, \cdots, \beta_p)$ contains (fixed) mean parameters, vector ${\bf u}_i$ contains random effects 

for the mean function, and 

 ${\bf e}_i = (e_{i1}, e_{i2}, \cdots, e_{in_i})^T$ are within-individual errors. 

We assume that ${\bf u}_i \sim N(0, D)$, and $e_{i1}, \cdots e_{in_i} $ are conditionally independent 

given ${\bf u}_i$, with 

$ e_{ij} \sim N(0, \sigma_{ij}^2)$. Here we allow the variances of the within-individual random 

errors, $\sigma^2_{ij}$, to be individual-specific and time-dependent. This allows great 

flexibility in modeling complicated longitudinal data, as discussed in the previous section. 

Let $\bsigma_i = (\sigma_{i1}^2, \cdots, \sigma_{in_i}^2)$. 



   To understand the nature of the within-individual variations and check if there are 

some systematic components in these variations, we investigate if 

the within-individual variances $\sigma^2_{ij}$ 

may be partially explained by time-varying covariates (say) ${\bf z}_i = (z_{i1}, \cdots, 

z_{in_i})$, where $z_{ij}$ is the value of covariate $z$ for individual $i$ at time $t_{ij}$. 

 Note that covariate ${\bf z}_i$ may be different from covariates ${\bf x}_i$. 

For simplicity, here we only consider one covariate, but the method can be easily extended 

to more than one covariates. 

Moreover, the within-individual variations $\sigma^2_{ij}$ may vary substantially across individuals. 

Thus, we model the within-individual variations and introduce additional random effects to 

account for the associations within an individual and variations across individuals, i.e., 

we consider the following model for the variance 

\begin{eqnarray}

     \log(\sigma_{ij}^2) = \alpha_1 + \alpha_2 z_{ij} + \eta a_i,  

   \qquad i=1,\cdots,n,  \quad j=1,2,\cdots, n_i, \label{var_model}

\end{eqnarray}

where vector $\balpha=(\alpha_1, \alpha_2)$ and $\eta$ are fixed parameters, 

and $a_i$ is a random effect for the variance model. 



We may assume that $\exp(a_i)$ follows 

an inverse gamma distribution or $a_i \sim N(0, 1)$. 

In the variance model (\ref{var_model}), we model the log-transformation of the 

variances, $\log(\sigma_{ij}^2)$, 

since the variances $\sigma_{ij}^2$ are strictly positive and the range of 

the log-transformation $\log(\sigma_{ij}^2)$ is 

$(-\infty, \infty)$, which allows us to introduce arbitrary covariate in the model. 



  In the above mean and variance models, we jointly model the conditional mean

$E({\bf y}_i | {\bf u}_i) = g({\bf x}_i, \bbeta, {\bf u}_i)$ and the conditional variance

$Var({\bf y}_i | {\bf u}_i) = diag(\sigma_{i1}^2, \cdots, \sigma_{in_i}^2)$. Modelling the

variances allows for more efficient estimation of the mean parameters and better understanding

of the within-individual variabilities.



\subsection{Robust model}



    In the variance model (\ref{var_model}), if we assume that 

$ \exp(a_i) \sim k/\chi^2_k$, i.e., the inverse $\chi^2$ distribution, 

which is a special case of 

a inverse gamma distribution, 

  then the random errors $e_{ij}$ for the within-individual repeated measurement 

follow  a $t(k)$-distribution with degrees of freedom $k$ 

(Lang et al. 1989). The $t(k)$-distribution has heavier tails than the normal distribution,

 and thus it incorporates possible 

outliers in the within-individual repeated measurements. 

The smaller the degrees of freedom $k$, the heavier the tails. 

The $t(k)$-distribution reduces to the standard 

 normal distribution when $k \to \infty$. Therefore, in this case models (\ref{nlme}) and 

(\ref{var_model}) may be used for robust inference. 



   Similarly, we may allow the variances for the random effects for the mean model ${\bf u}_i$ 

to be individual-specific, i.e., ${\bf u}_i \sim N(0, D_i)$, and model the covariance matrix 

$D_i$ in a similar way as $\sigma^2_{ij}$. While the approaches are conceptually similar, 

for simplicity we focus on the variance model (\ref{var_model}) and assume $D_i=D$ in this 

paper. 



\subsection{Covariate measurement errors}



   The covariate $z_{ij}$ in the variance model may be measured with errors, such as CD4 cell count. 

The covariate effect on the variance, if any, may be masked if the measurement errors are 

not addressed. 

To address the measurement errors, we may view the repeated measurements of 

the covariate as its ``replicates" in order to estimate the magnitude of the measurement errors

and model the covariate process using an empirical LME model. 



   Let $z^*_{ij}$ be the unobserved

true covariate value for individual $i$ at time $t_{ij}$ and let $z_{ij}$ be the corresponding

observed value with measurement error. We 

may consider the following mixed effects model as a classic measurement error model 

\begin{eqnarray}

{z}_{ij} \equiv {z}^*_{ij} + \epsilon_{ij} = q(t_{ij}, \bgamma, {\bf b}_i) + \epsilon_{ij}, 

\label{CD4}

\end{eqnarray}

where $q(\cdot)$ is a known linear function or a unknown smooth function, $\bgamma$ contains 

fixed parameters, ${\bf b}_i$ contains random effects, and 

$\epsilon_{ij}$ is measurement error. 

We assume that $b_i \sim N(0, B),$ and independently $ \epsilon_{ij} \; i.i.d. \; \sim N(0,

\omega^2)$.



For example, we may model the CD4 process using 

a quadratic mixed effects model ${z}_{ij} = \gamma_0+\gamma_1 t_{ij} + \gamma_2 t_{ij}^2

     + b_i + \epsilon_{ij}$. Alternatively, for more complex CD4 trajectories, we 

may choose $q(t_{ij}, \bgamma, {\bf b}_i)$ to be a nonparametric smooth function, then 

approximate it using a basis-based approach, which leads to an approximate LME model (Wu, 2009). 

In either case, we may consider an LME model for the  measurement error model (\ref{CD4}). 



\section{Approximate Likelihood Inference}



    Let $\btheta = (\bbeta, \bgamma, \bxi)$ denote 

the collection of all fixed parameters in the models, 

where $(\bbeta, \bgamma)$ contains fixed mean parameters, and $\bxi = (\balpha, D, \eta, 

\omega^2)$ 

contains fixed variance (dispersion) parameters. 

Let ${\bf v}_i = ({\bf u}_i, a_i, {\bf b}_i)$ denote the collection of all random effects, 

and let ${\bf v} = ({\bf v}_1, \cdots, {\bf v}_n)$. 

Standard likelihood inference for $\btheta$ is based on the following 

(marginal) likelihood function, which integrates out the unobservable random effects  

\begin{eqnarray*}

 L(\btheta) = \prod_{i=1}^n \int\int\int f({\bf y}_i | \; {\bf x}_i, {\bf z}_i, 

  \bbeta, D, \bsigma^2_{i}, {\bf u}_i)

  f(\bsigma^2_{i} | \balpha, a_i) f({\bf z}_i | \bgamma, {\bf b}_i) f({\bf u}_i, {\bf b}_i, a_i |

\btheta) \; d {\bf u}_i \; da_i \; d{\bf b}_i,     

 %% = \prod_{i=1}^m \int f({\bf y}_i | \btheta, {\bf v}_i) 

 %% f({\bf v}_i | \btheta) \; d {\bf v}_i. 

\end{eqnarray*} 

where $f(\cdot)$ denotes a generic density function. 

The above likelihood involves a high-dimensional and 

intractable integration.

Monte Carlo EM algorithms or Gaussian quadrature numerical 

integration methods can be computationally very intensive. Thus, in the following, we discuss 

a computationally more efficient  approximate method  for approximate likelihood inference, using  

the h-likelihood (Lee et al. 2017) or conditional second-order generalized estimating equations 

(CGEE2, Vonesh et al. 2002). 





   The {h-loglikelihood} (Lee et al. 2017) can be written as 

$$ h(\btheta, {\bf v}) = \sum_{i=1}^n \left\{\log f({\bf y}_i | \; {\bf x}_i, {\bf z}_i,

  \bbeta, D, \bsigma^2_{i}, {\bf u}_i) + 

  \log f(\bsigma^2_{i} | \balpha, a_i) + \log f({\bf z}_i | \bgamma, {\bf b}_i) 

  + \log f({\bf u}_i, {\bf b}_i, a_i |\btheta)\right\}  $$

where the random effects ${\bf v}_i = ({\bf u}_i, {\bf b}_i, a_i)$ are treated as ``parameters". 

Given parameter $\btheta$, estimates of the random effects ${\bf v}$ can be obtained by maximizing the 

h-loglikelihood $h(\btheta, {\bf v})$ with respect to ${\bf v}$, denoted by $\hat{\bf v}$, 

i.e., $\hat{\bf v}$ solves

$\partial h(\btheta, {\bf v})/\partial{\bf v}=0$. 

Then, we can obtain estimates of the parameters $\btheta$ by maximizing $h(\btheta, \hat{\bf v})$. 

The process iterates until convergence. 

For binary or discrete response ${\bf y}_i$, Lee et al. (2017) recommends adding an adjusted 

term to reduce bias and obtain estimates of the parameters $\btheta$ by maximizing the following 

``adjusted profile likelihood"

  $$   \log L(\btheta)  \approx p_{\bf v}(h)

         \equiv \left[h(\btheta, {\bf v})

         - \frac{1}{2}\log det\{H(h; {\bf v})/(2\pi)\}\right] \Bigg |_{{\bf v}=\hat{\bf v}},   $$

where the second term is an adjusted term, and 

$H(h; {\bf v}) = - \partial^2 h/\partial{\bf v}^2$ and $\hat{\bf v}$ solves

$\partial h/\partial{\bf v}=0$. Lee and Nelder (2001) show that the h-likelihood estimate 

is the first-order Laplace approximation to the marginal log-likelihood 

$\log L(\btheta)$.  





     Vonesh et al. (2002) proposed the CGEE2 method which also maximizes $h(\btheta, {\bf v})$

but without the adjusted term. That is, the CGEE2 estimate may be viewed as the 

h-likelihood estimate without the adjusted term. 

They showed that their CGEE2 estimate is asymptotically equivalent 

to the h-likelihood estimate, the linearization estimate of 

Lindstrom and Bates (1990) for NLME models, and the penalized quasi-likelihood (PQL) estimate of 

Breslow and Clayton (1993), when both the sample size $n$ and the numbers of repeated 

measurements $\min_i\{n_i\}$ go to $\infty$. Moreover, they showed that 

the CGEE2 estimates are consistent and asymptotically normally distributed. 

Following their ideas, we propose the following iterative algorithm to obtain approximate 

MLE of the parameter $\btheta = (\bbeta, \bgamma, \bxi)$. Starting with some initial values 

$(\btheta^{(0)}, {\bf v}^{(0)})$, at iterative $k, k = 0,1,2,\cdots,$ 

\begin{itemize}

\item Step 1: Solve equation $\partial h({\btheta}^{(k)}, {\bf v})/\partial{\bf v}=0$ and obtain 

updated estimate of the random effects $\hat{\bf v}^{(k+1)}$; 

\item Step 2: Solve equation 

$\partial h(\bbeta, {\bgamma}, {\bxi}^{(k)}; {\bf v}^{(k+1)})/\partial({\bbeta},{\bgamma})=0$ 

and obtain  

updated estimates of the mean parameters $(\bbeta^{(k+1)},{\bgamma}^{(k+1)})$; 

\item  Step 3: Solve equation

$\partial h(\bbeta^{(k+1)}, {\bgamma}^{(k+1)}, {\bxi}; {\bf v}^{(k+1)})/\partial{\bxi}=0$ and obtain

updated estimates of the variance-covariance parameters $\bxi^{(k+1)}$.

\end{itemize}

Iterate the above steps until convergence, we obtain approximate MLE $\hat{\btheta}$. 

Standard errors of the estimates may be obtained from the diagonal elements (square roots) of 

the inverse of the Fisher information matrix. 





\section{A Real Data Example}



    Analyze a real HIV/AIDS dataset using three methods: the Linstrom and Bates (1990) method 

implemented in R package {\em nlme} without modeling the variances, the two-step method, and 

the proposed method. The results should show efficiency gain (smaller standard errors) 

of the estimates $\hat{\bbeta}$ which allow us to compare different treatments more 

effectively. Estimates of $\balpha$ also allow us to see if CD4 partially explains the 

within-individual variance $\sigma_{ij}^2$. 





\section{Simulation Study}



    In the simulation study, we will evaluate the proposed method and compare it with 

the Linstrom and Bates (1990) method implemented in R package {\em nlme} without modeling the

variances and the two-step method. The simulation will be designed in the following ways: 

i) compare estimates, standard errors, MSEs, and coverage probabilities; ii) bootstrap 

standard errors (either parametric or nonparametric bootstrap); iii) robustness of the 

proposed method (with inverse-gamma distribution for $a_i$) for simulated data with 

5\% outliers (adding a large number, say 10, to 5\% simulated data); iv) different 

sample size $n$ (say $n=100, 200$) and different values of $n_i$ (say, 15, 30). 



\section{Discussion}



\section*{References}



\begin{description}

\item German, C.A., Sinsheimer, J.S., Zhou, J., and Zhou, H. (2022). 

  WiSER: Robust and scalable estimation and inference of

within-subject variances from intensive longitudinal data. {\it Biometrics}, in press. 



\item Hedeker, D., Mermelstein, R.J., and Demirtas, H. (2008). 

  An Application of a Mixed-Effects Location Scale Model for Analysis

of Ecological Momentary Assessment (EMA) Data. {\it Biometrics 64}, 627-634. 



\item Lee, Y.,  Nelder, J.A., and Pawitan, Y. (2017). 

{\it Generalized Linear Models with Random Effects: Unified Analysis via 

H-likelihood}, Second edition. London: Chapman \& Hall/CRC.  



\item Lin, X., Raz, J., and Harlow, S.D. (1997). 

  Linear Mixed Models with Heterogeneous Within-Cluster Variances. {\it Biometrics 53}, 910-923. 



\item Lindstrom, M.J. and Bates, D.M. (1990).   Nonlinear  mixed

       effects models for repeated measures data. {\it Biometrics {\bf 46}},

       673-687. 



 \item Noh, M., Wu, L., and Lee, Y. (2012).

        Hierarchical Likelihood Methods for Nonlinear and

        Generalized Linear Mixed Models with Missing Data and

        Measurement Errors in Covariates, {\it Journal of Multivariate Analysis},

        109, 42-51. 



\item Pourahmadi, M. (1999). Joint mean-covariance models with applications to longitudinal data:

Unconstrained parameterisation. {\it Biometrika 86}, 677-690. 



\item  Vonesh, E.F., Wang, H., Nie, L., and Majumdar, D. (2002). 

  Conditional second-order generalized estimating equations for generalized 

  linear and nonlinear mixed-effects models. 

 {\it Journal of the American Statistical Association {\bf 97}}, 271-283.





\end{description}



\end{document}



